\subsection{Background: Systemic Tau and Antisynchrony}

Systemic Tau $\tau_s(k)$ is a local measure of ordinal persistence derived from permutation entropy over sliding windows. Antisynchrony is defined as
\[
A(k) = \mathrm{std}\bigl(\{\tau_s^i(k)\}_{i=1}^N\bigr),
\]
where the standard deviation is taken across the $N$ modules (or regions) at time $k$. Low sustained values of $A(k)$ indicate alignment of persistence scales (Layer 2).

\subsection{Layer 1 Intensification}

Layer 1 is identified by excursions of the critical mass metric $M(k)$, which combines normalized hyper-persistence of $\tau_s$ with RQA-derived trapping time (or laminarity) under Config B parameters.

\subsection{Definition of the Joint Episode}

\begin{conceptbox}[Definition: Joint Episode]
A \textbf{Joint Episode} is a maximal contiguous interval $[t_\mathrm{start}, t_\mathrm{end}]$ such that:
\begin{enumerate}
    \item $A(k) < \theta_A$ for all $k \in [t_\mathrm{start}, t_\mathrm{end}]$ and $t_\mathrm{end} - t_\mathrm{start} + 1 \ge D_{\min}$ (sustained Layer 2),
    \item there exists at least one $k^* \in [t_\mathrm{start}, t_\mathrm{end}]$ with $M(k^*) \ge \theta_M$ (presence of Layer 1 intensification).
\end{enumerate}
The interval is maximal: it cannot be extended without violating either condition.
\end{conceptbox}

The thresholds $\theta_A$, $D_{\min}$, and $\theta_M$ are regime-calibrated (see Section 3).

\begin{figure}[htbp]
\centering
\begin{tikzpicture}[
  node distance=0.55cm and 0.4cm,
  layer box/.style={
    draw=conceptcolor!65, 
    fill=conceptcolor!7, 
    thick, 
    rounded corners=2.5pt, 
    minimum width=10.2cm, 
    minimum height=1.05cm, 
    align=center, 
    font=\small\bfseries,
    inner sep=4pt
  },
  sublabel/.style={font=\footnotesize, align=center},
  je box/.style={
    draw=orange!75!black, 
    fill=orange!10, 
    thick, 
    rounded corners=3pt, 
    minimum width=8.8cm, 
    minimum height=0.85cm, 
    align=center, 
    font=\small\bfseries,
    inner sep=3pt
  },
  arrow up/.style={-{Stealth[length=2.8mm]}, very thick, conceptcolor!75},
  arrow down/.style={-{Stealth[length=2.8mm]}, very thick, gray!65, dashed},
  time arrow/.style={-{Stealth}, thick, black!70}
]

% === Vertical stack: L3 (top) to L1 (bottom) ===
\node[layer box] (L3) at (0, 3.6) {
  Layer 3: Structural Reorganization \\
  \sublabel{Global topological or dynamical reconfiguration of the system}
};

\node[layer box] (L2) at (0, 2.15) {
  Layer 2: Relational Coherence / Permeability \\
  \sublabel{Sustained low antisynchrony $A(k)$; alignment of persistence scales across modules}
};

\node[layer box] (L1) at (0, 0.7) {
  Layer 1: Local Intensification of Persistence \\
  \sublabel{Systemic Tau $\tau_s(k)$ peaks + critical-mass metric $M(k) \ge \theta_M$}
};

% === Joint Episode box (spans L1--L2) ===
\node[je box] (JE) at (0, 1.42) {
  \textbf{Joint Episode} (kairos) \\
  \sublabel{maximal interval of low $A(k)$ containing Layer-1 intensification}
};

% === Time axis ===
\draw[time arrow] (-5.8, -0.35) -- (5.8, -0.35) 
  node[right, font=\footnotesize, black!75] {discrete time $k$};

% Episode duration marker on time axis
\draw[thick, orange!80!black, line width=1.1pt] (-3.6, -0.35) -- (2.4, -0.35);
\node[font=\tiny, orange!85!black, below=1.5pt] at (-0.6, -0.35) {$\,[t_\mathrm{start},\, t_\mathrm{end}]\,$};

% Vertical guides from JE to time (subtle)
\draw[orange!40, thin, dashed] (-3.6, -0.1) -- (-3.6, -0.28);
\draw[orange!40, thin, dashed] (2.4, -0.1) -- (2.4, -0.28);

% === Bottom-up enabling arrows (left side) ===
\draw[arrow up] ([xshift=-3.8cm]L1.west) -- ++(0,0.55) -| 
  node[near start, left=1pt, font=\footnotesize, text=conceptcolor!80] {bottom-up enabling} 
  ([xshift=-3.8cm]L2.west);
\draw[arrow up] ([xshift=-4.3cm]L2.west) -- ++(0,0.55) -| 
  ([xshift=-4.3cm]L3.west);

% === Top-down constraint arrows (right side) ===
\draw[arrow down] ([xshift=4.3cm]L3.east) -- ++(0,-0.55) -| 
  node[near start, right=1pt, font=\footnotesize, text=gray!70] {top-down constraints} 
  ([xshift=4.3cm]L2.east);
\draw[arrow down] ([xshift=3.8cm]L2.east) -- ++(0,-0.55) -| 
  ([xshift=3.8cm]L1.east);

% === Small annotation labels inside layers ===
\node[font=\tiny, conceptcolor!60, right=2pt] at (L1.east) {$\tau_s$, $M$};
\node[font=\tiny, conceptcolor!60, right=2pt] at (L2.east) {$A(k)$};
\node[font=\tiny, gray!55, left=2pt] at (L3.west) {reorg.};

% === Legend / key (bottom right) ===
\node[draw=gray!40, fill=white, rounded corners=1.5pt, 
      font=\footnotesize, inner sep=2pt, anchor=south east] 
  at (5.6, -0.9) { 
  \textcolor{conceptcolor!75}{$\uparrow$} enabling \quad 
  \textcolor{gray!65}{$\downarrow$} constraint 
};

\end{tikzpicture}
\caption{Three-layer ontology and the Joint Episode as \textit{kairos}. Layer 1 (bottom) represents local intensification of persistence through peaks in Systemic Tau $\tau_s(k)$ and the critical-mass metric $M(k)$. Layer 2 (middle) captures relational coherence and permeability via sustained low antisynchrony $A(k)$ across modules. The Joint Episode is defined as their joint realization: a maximal interval of low $A(k)$ that contains at least one Layer-1 intensification. Upward arrows (left) indicate bottom-up enabling from local persistence to scale coherence to structural possibility; downward arrows (right) indicate top-down constraints whereby Layer-3 organization restricts admissible L1--L2 configurations. The highlighted interval on the time axis marks the structured window of the present.}
\label{fig:three_layers_je}
\end{figure}

